\section{Conclusion}

RQ1 asks which streaming communication technologies, frameworks, and interaction patterns are suitable for large-scale and continuous data exchange in microservice-based systems such as DYNAMOS. At the level of isolated transport performance, client-streaming gRPC is the strongest raw result in the generic pipeline benchmark. Batched unary is close in the clean matrix and stronger in the backpressure scenario, while per-message unary is much slower for bulk transfer. That result establishes the value of amortising per-message overhead, whether by streaming or by batching, rather than treating streaming itself as the only useful mechanism. The brokered alternatives expose the cost and benefit of decoupling, persistence, replay, and failure handling.

RQ2 asks how streaming communication can be integrated into DYNAMOS while preserving dynamic service composition, policy-driven orchestration, and isolation. The implementation keeps the existing control path intact: requests are still evaluated and composed before execution, and generated jobs remain the execution unit. The change is in the bulk data path. Provider-side chunks must remain chunks through aggregate, algorithm, agent, and gateway stages; otherwise the system merely hides a full-buffered result behind a streaming interface.

RQ3 validates the implemented DYNAMOS variants on local bulk \texttt{dataThroughTtp} retrieval. All chunk-aware variants beat the classic unary baseline in average completion time and peak memory for 50k, 100k, and 250k row result sets. They also expose earlier partial results, which classic unary cannot do because its first result is its final result. Focused checks on \texttt{computeToData} and multi-provider \texttt{dataThroughTtp} show that the implementation is not confined to one request shape. The strongest local recommendation is chunked unary. It gives the best 100k and 250k averages, the lowest 250k memory peak, and the smallest architectural shift from the original message-driven design.

The more general conclusion is boundary-specific. Intra-job gRPC streaming is useful where a direct stream between generated services is acceptable. RabbitMQ Streams is not the source of the local speedup by itself, because that speedup mainly comes from chunk preservation. Its stronger case is distributed inter-agent transfer, where retained records, replay, and independent consumers may matter more than local latency. The Scattered Directive and FABRIC context is therefore future work: it moves DYNAMOS toward cross-node VFL workflows, where network separation, broker placement, policy changes, and multi-party scaling can dominate local transfer costs.

The main research question asks how streaming communication can be integrated into DYNAMOS without weakening its policy-driven and adaptive character. The answer is to treat streaming as a data-path change after policy approval, not as a replacement for policy enforcement or orchestration. Compliance is preserved by construction in the evaluated implementation because the policy decision point, generated job composition, provider selection, and data-steward isolation remain unchanged. What changes is how approved bulk results are transferred once a job is already allowed to run.

